\chapter{Επίδειξη Ελέγχου Πρόσβασης}
\label{ch:demonstration}

Το παρόν κεφάλαιο παρουσιάζει πειραματικά τον κύκλο εκχώρησης πρόσβασης
του \textlatin{Plegma Data Space} μέσω σεναρίων επίδειξης. Κάθε σενάριο
εκτελείται σε πραγματικές συνθήκες λειτουργίας και επαληθεύει μια διαφορετική
πτυχή της αρχιτεκτονικής ελέγχου πρόσβασης. Τα σενάρια οργανώνονται σε δύο
επίπεδα: η πρώτη ομάδα αφορά τον έλεγχο πρόσβασης του παρόχου, ενώ η δεύτερη
τον κύκλο εκχώρησης εξουσιοδότησης στον καταναλωτή.

% ------------------------------------------------------------------
\section{Περιβάλλον Επίδειξης}
\label{sec:demo-environment}
% ------------------------------------------------------------------

Η επίδειξη εκτελείται σε τοπικό περιβάλλον \textlatin{Windows} με τα
\textlatin{containers} να εκτελούνται μέσω \textlatin{Docker Desktop}. Το
\textlatin{Kong API Gateway} εκτίθεται στη θύρα \textlatin{8000} του τοπικού
μηχανήματος και αποτελεί το μοναδικό σημείο εισόδου για όλα τα αιτήματα
δεδομένων. Ο \textlatin{Keyrock} εκτίθεται στη θύρα \textlatin{3007} για τη
διαδικασία ανταλλαγής \textlatin{token}.

Τα εργαλεία που χρησιμοποιούνται για την επίδειξη είναι:
\begin{itemize}
    \item \textlatin{PowerShell}: Για την εκτέλεση των εντολών
    επίδειξης.
    \item \textlatin{curl}: Για την αποστολή \textlatin{HTTP}
    αιτημάτων.
    \item \textlatin{Docker}: Για την εκτέλεση των \textlatin{scripts}
    δημιουργίας \textlatin{JWT}.
    \item \textlatin{Python}: Για την παρουσίαση των αποτελεσμάτων
    σε \textlatin{JSON} μορφή (\textlatin{python -m json.tool}).
\end{itemize}

% ------------------------------------------------------------------
\section{Σενάριο 1: Πρόσβαση Χωρίς Ταυτοποίηση}
\label{sec:scenario1}
% ------------------------------------------------------------------

Το πρώτο σενάριο επαληθεύει ότι το \textlatin{Kong API Gateway} απορρίπτει
αιτήματα που δεν φέρουν διαπιστευτήρια ταυτοποίησης. Αυτό αποτελεί το
θεμελιώδες χαρακτηριστικό ελέγχου πρόσβασης: κανένα αίτημα δεν μπορεί να
φτάσει στα δεδομένα χωρίς έγκυρο \textlatin{token}.

\selectlanguage{english}
\begin{lstlisting}[language=bash, caption={Request without authentication token.}, stringstyle=\color{black}, keywordstyle=\color{black}]
curl http://localhost:8000/ngsi-ld/v1/entities \
     ?type=https://plegma.example.org/vocab%23HouseholdElectricMeasurement \
     &limit=1
\end{lstlisting}
\selectlanguage{greek}

\noindent Αναμενόμενο αποτέλεσμα:
\selectlanguage{english}
\begin{lstlisting}[language=bash, caption={Kong response for unauthenticated request.}, stringstyle=\color{black}, keywordstyle=\color{black}]
HTTP/1.1 401 Unauthorized
{"message": "Unauthorized"}
\end{lstlisting}
\selectlanguage{greek}
Το \textlatin{Kong} επιστρέφει \textlatin{HTTP 401 Unauthorized} καθώς
δεν υπάρχει κεφαλίδα \textlatin{Authorization} στο αίτημα. Το αίτημα
δεν προωθείται ποτέ στον \textlatin{Orion-LD}.

% ------------------------------------------------------------------
\section{Σενάριο 2: Πρόσβαση με \textlatin{Token} Παρόχου}
\label{sec:scenario2}
% ------------------------------------------------------------------

Το δεύτερο σενάριο επαληθεύει τη ροή ανάγνωσης δεδομένων από τον πάροχο.
Δημιουργείται ένα \textlatin{iSHARE JWT} με το ιδιωτικό κλειδί του παρόχου
και χρησιμοποιείται για πρόσβαση στα δεδομένα του \textlatin{Plegma Dataset}.

\subsection{Δημιουργία \textlatin{iSHARE JWT} Παρόχου}
\label{subsec:scenario2-token}

\selectlanguage{english}
\begin{lstlisting}[language=bash, caption={iSHARE JWT generation via Docker.}, stringstyle=\color{black}, keywordstyle=\color{black}]
PROVIDER_TOKEN=$(docker run --rm \
  -v "./iShare:/iShare" \
  python:3.9-slim sh -c \
  "pip install PyJWT cryptography -q 2>/dev/null && \
   python3 /iShare/gen_token.py")
\end{lstlisting}
\selectlanguage{greek}
\noindent Τα ορίσματα ερμηνεύονται ως εξής:
\begin{itemize}
    \item \textlatin{-{}-rm}: Αφαιρεί το \textlatin{container} μετά την εκτέλεση.
    \item \textlatin{-v "./iShare:/iShare"}: Προσαρτά τον φάκελο με τα
    πιστοποιητικά στο \textlatin{container}.
    \item \textlatin{pip install PyJWT cryptography -q 2>/dev/null}: Εγκαθιστά
    τις απαραίτητες βιβλιοθήκες αθόρυβα.
    \item \textlatin{python3 /iShare/gen\_token.py}: Εκτελεί το \textlatin{script}
    δημιουργίας \textlatin{JWT}.
\end{itemize}

\subsection{Ανάγνωση Δεδομένων \textlatin{HouseholdElectricMeasurement}}
\label{subsec:scenario2-electric}

\selectlanguage{english}
\begin{lstlisting}[language=bash, caption={Electric measurement read request with provider token.}, stringstyle=\color{black}, keywordstyle=\color{black}]
curl http://localhost:8000/ngsi-ld/v1/entities \
     ?type=https://plegma.example.org/vocab%23HouseholdElectricMeasurement \
     &limit=1 \
     -H "Authorization: Bearer $PROVIDER_TOKEN"
\end{lstlisting}
\selectlanguage{greek}
\noindent \textbf{Αποτέλεσμα (\textlatin{HTTP 200 OK}):}
\selectlanguage{english}
\begin{lstlisting}[language=bash, caption={Response with electric measurement data.}, stringstyle=\color{black}, keywordstyle=\color{black}]
[
  {
    "id": "urn:ngsi-ld:HouseholdElectricMeasurement:House01:2022-07-15T00:00:00Z",
    "type": "https://plegma.example.org/vocab#HouseholdElectricMeasurement",
    "https://smartdatamodels.org/dateObserved": {
      "type": "Property",
      "value": { "@type": "DateTime", "@value": "2022-07-15T00:00:00Z" }
    },
    "https://plegma.example.org/vocab#activePower": {
      "type": "Property",
      "value": 121.542,
      "unitCode": "WTT"
    },
    "https://plegma.example.org/vocab#phaseVoltage": {
      "type": "Property",
      "value": 236.83,
      "unitCode": "VLT"
    },
    "https://plegma.example.org/vocab#appliancePower": {
      "type": "Property",
      "value": { "ac_1": 0, "ac_2": 2.033, "boiler": 0 },
      "unitCode": "WTT"
    }
  }
]
\end{lstlisting}
\selectlanguage{greek}
\subsection{Ανάγνωση Περιβαλλοντικών Μετρήσεων}
\label{subsec:scenario2-env}

Ο ίδιος \textlatin{token} παρόχου επιτρέπει πρόσβαση και στους υπόλοιπους
τύπους οντοτήτων που διαθέτει το \textlatin{Plegma Dataset}:

\selectlanguage{english}
\begin{lstlisting}[language=bash, caption={Environmental measurement read request.}, stringstyle=\color{black}, keywordstyle=\color{black}]
curl http://localhost:8000/ngsi-ld/v1/entities \
     ?type=https://plegma.example.org/vocab%23EnvironmentalMeasurement \
     &limit=1 \
     -H "Authorization: Bearer $PROVIDER_TOKEN"
\end{lstlisting}
\selectlanguage{greek}
\noindent \textbf{Αποτέλεσμα (\textlatin{HTTP 200 OK}):}
\selectlanguage{english}
\begin{lstlisting}[language=bash, caption={Response with environmental measurement data.}, stringstyle=\color{black}, keywordstyle=\color{black}]
[
  {
    "id": "urn:ngsi-ld:EnvironmentalMeasurement:House01:2022-07-15T00:00:00Z",
    "type": "https://plegma.example.org/vocab#EnvironmentalMeasurement",
    "https://plegma.example.org/vocab#indoorTemperature": {
      "type": "Property",
      "value": 27.4,
      "unitCode": "CEL"
    },
    "https://plegma.example.org/vocab#indoorHumidity": {
      "type": "Property",
      "value": 32,
      "unitCode": "P1"
    }
  }
]
\end{lstlisting}
\selectlanguage{greek}
\subsection{Ανάγνωση Δεδομένων Κτιρίου}
\label{subsec:scenario2-building}

\selectlanguage{english}
\begin{lstlisting}[language=bash, caption={Building data read request.}, stringstyle=\color{black}, keywordstyle=\color{black}]
curl http://localhost:8000/ngsi-ld/v1/entities \
     ?type=https://smartdatamodels.org/dataModel.Building/Building \
     &limit=1 \
     -H "Authorization: Bearer $PROVIDER_TOKEN"
\end{lstlisting}
\selectlanguage{greek}
\noindent \textbf{Αποτέλεσμα (\textlatin{HTTP 200 OK}):}
\selectlanguage{english}
\begin{lstlisting}[language=bash, caption={Response with building metadata.}, stringstyle=\color{black}, keywordstyle=\color{black}]
[
  {
    "id": "urn:ngsi-ld:Building:House01",
    "type": "https://smartdatamodels.org/dataModel.Building/Building",
    "https://plegma.example.org/vocab#heatingSystem": {
      "type": "Property",
      "value": ["radiatorOil", "airConditioner"]
    },
    "https://plegma.example.org/vocab#hasAppliance": {
      "type": "Relationship",
      "object": [
        "urn:ngsi-ld:Device:House01:ac_1",
        "urn:ngsi-ld:Device:House01:ac_2",
        "urn:ngsi-ld:Device:House01:boiler"
      ]
    },
    "numberOfRooms": { "type": "Property", "value": 3 }
  }
]
\end{lstlisting}
\selectlanguage{greek}
Και στις τρεις περιπτώσεις, το \textlatin{Kong} επαλήθευσε επιτυχώς το
\textlatin{iSHARE JWT} και προώθησε τα αιτήματα στον \textlatin{Orion-LD},
ο οποίος επέστρεψε τα αντίστοιχα δεδομένα με \textlatin{HTTP 200 OK}.

% ------------------------------------------------------------------
\section{Σενάριο 3: Κύκλος Εκχώρησης Πρόσβασης στον Καταναλωτή}
\label{sec:scenario3}
% ------------------------------------------------------------------

Το τρίτο σενάριο αποτελεί την πλήρη επίδειξη του κύκλου εκχώρησης
εξουσιοδότησης (\textlatin{delegation cycle}): ο πάροχος έχει εκχωρήσει άδεια
στον καταναλωτή να διαβάζει δεδομένα τύπου \textlatin{HouseholdElectricMeasurement},
και η επίδειξη αυτή αποδεικνύει ότι ο καταναλωτής μπορεί να αποκτήσει
πρόσβαση αποκλειστικά και μόνο μέσω αυτής της εκχώρησης.

\subsection{Βήμα 1: Δημιουργία \textlatin{iSHARE JWT} Καταναλωτή}
\label{subsec:scenario3-consumer-jwt}

\selectlanguage{english}
\begin{lstlisting}[language=bash, caption={Consumer iSHARE JWT generation.}, stringstyle=\color{black}, keywordstyle=\color{black}]
CONSUMER_JWT=$(docker run --rm \
  -v "./iShare:/iShare" \
  python:3.9-slim sh -c \
  "pip install PyJWT cryptography -q 2>/dev/null && \
   python3 /iShare/gen_consumer_token.py")
\end{lstlisting}
\selectlanguage{greek}
\subsection{Βήμα 2: Λήψη \textlatin{Access Token} από τον \textlatin{Keyrock}}
\label{subsec:scenario3-access-token}

\selectlanguage{english}
\begin{lstlisting}[language=bash, caption={iSHARE JWT exchange for access token.}, stringstyle=\color{black}, keywordstyle=\color{black}]
TOKEN_RESPONSE=$(curl -s -X POST \
  "http://localhost:3007/oauth2/token" \
  -H "Content-Type: application/x-www-form-urlencoded" \
  -d "grant_type=client_credentials \
      &scope=iSHARE \
      &client_assertion_type=urn:ietf:params:oauth:client-assertion-type:jwt-bearer \
      &client_assertion=${CONSUMER_JWT} \
      &client_id=EU.EORI.NLPLEGMACONSUMER")

ACCESS_TOKEN=$(echo $TOKEN_RESPONSE | python -c \
  "import sys, json; print(json.load(sys.stdin)['access_token'])")
\end{lstlisting}
\selectlanguage{greek}
Η επιτυχής λήψη \textlatin{access token} από τον \textlatin{Keyrock}
υποδηλώνει ότι:
\begin{enumerate}
    \item Το \textlatin{iSHARE JWT} του καταναλωτή είναι έγκυρο.
    \item Ο καταναλωτής είναι εγγεγραμμένος ενεργός συμμετέχων στον
    \textlatin{Satellite}.
    \item Υπάρχει ενεργή πολιτική εκχώρησης εξουσιοδότησης από τον πάροχο
    προς τον καταναλωτή.
\end{enumerate}

\subsection{Βήμα 3: Πρόσβαση στα Δεδομένα Μέσω \textlatin{Kong}}
\label{subsec:scenario3-data-access}

\selectlanguage{english}
\begin{lstlisting}[language=bash, caption={Data read request by the consumer.}, stringstyle=\color{black}, keywordstyle=\color{black}]
curl http://localhost:8000/ngsi-ld/v1/entities \
     ?type=https://plegma.example.org/vocab%23HouseholdElectricMeasurement \
     &limit=1 \
     -H "Authorization: Bearer $ACCESS_TOKEN"
\end{lstlisting}
\selectlanguage{greek}
\noindent \textbf{Αποτέλεσμα (\textlatin{HTTP 200 OK}):}
\selectlanguage{english}
\begin{lstlisting}[language=bash, caption={Successful data read by the consumer.}, stringstyle=\color{black}, keywordstyle=\color{black}]
[
  {
    "id": "urn:ngsi-ld:HouseholdElectricMeasurement:House01:2022-07-15T00:00:00Z",
    "type": "https://plegma.example.org/vocab#HouseholdElectricMeasurement",
    "https://plegma.example.org/vocab#activePower": {
      "type": "Property",
      "value": 121.542,
      "unitCode": "WTT"
    },
    "https://plegma.example.org/vocab#phaseVoltage": {
      "type": "Property",
      "value": 236.83,
      "unitCode": "VLT"
    }
  }
]
\end{lstlisting}
\selectlanguage{greek}
Το \textlatin{Kong} επαλήθευσε το \textlatin{access token}, εξέτασε το
\textlatin{delegationEvidence} που ενσωματώνει, διαπίστωσε ότι ο τύπος
πόρου και η ενέργεια \textlatin{GET} καλύπτονται από την πολιτική του
παρόχου, και προώθησε το αίτημα στον \textlatin{Orion-LD}.

% ------------------------------------------------------------------
\section{Σύνοψη Αποτελεσμάτων Επίδειξης}
\label{sec:demo-results}
% ------------------------------------------------------------------

Τα αποτελέσματα των τριών σεναρίων επίδειξης συνοψίζονται στον
Πίνακα~\ref{tab:demo-results}.

\begin{table}[H]
\centering
\caption{Αποτελέσματα σεναρίων επίδειξης ελέγχου πρόσβασης.}
\label{tab:demo-results}
\begin{tabular}{|>{\centering\arraybackslash}m{0.5cm}|>{\centering\arraybackslash}m{5cm}|>{\centering\arraybackslash}m{2.5cm}|>{\centering\arraybackslash}m{5cm}|}
\hline
\textbf{\#} & \textbf{Σενάριο} & \textbf{Αποτέλεσμα} & \textbf{Αιτιολογία} \\
\hline
1 & Χωρίς \textlatin{token} & \textlatin{401 Unauthorized} & Απουσία κεφαλίδας \textlatin{Authorization} \\
\hline
2α & Πάροχος, \textlatin{HouseholdElectricMeasurement} & \textlatin{200 OK} & Έγκυρο \textlatin{iSHARE JWT}, τοπική αρχή \\
\hline
2β & Πάροχος, \textlatin{EnvironmentalMeasurement} & \textlatin{200 OK} & Έγκυρο \textlatin{iSHARE JWT}, τοπική αρχή \\
\hline
2γ & Πάροχος, \textlatin{Building} & \textlatin{200 OK} & Έγκυρο \textlatin{iSHARE JWT}, τοπική αρχή \\
\hline
3 & Καταναλωτής, \textlatin{HouseholdElectricMeasurement} & \textlatin{200 OK} & Έγκυρο \textlatin{delegation evidence} \\
\hline
\end{tabular}
\end{table}

Η επίδειξη αποδεικνύει ότι το σύστημα λειτουργεί σύμφωνα με τις αρχές
ενός \textlatin{Data Space} βάσει \textlatin{iSHARE}:

\begin{itemize}
    \item Μηδενική πρόσβαση χωρίς ταυτοποίηση: Κανένα αίτημα
    χωρίς έγκυρο \textlatin{token} δεν φτάνει στα δεδομένα.
    \item Πλήρης πρόσβαση παρόχου: Ο πάροχος έχει πρόσβαση σε
    όλα τα δεδομένα που κατέχει, αποδεικνύοντας ταυτότητα μέσω
    ψηφιακής υπογραφής.
    \item Ελεγχόμενη πρόσβαση καταναλωτή: Ο καταναλωτής αποκτά
    πρόσβαση μόνο στα δεδομένα για τα οποία ο πάροχος έχει εκδώσει
    πολιτική εκχώρησης — διατηρώντας έτσι ο πάροχος την ψηφιακή του
    κυριαρχία.
\end{itemize}

